%% Pitts 1995 Figures 4, 6, 8: all fifteen directions of seven bi-rules.
\calculussheet{sheet:contextual-adjoint-first-order}{Pitts's adjoint contextual first-order rules}{Use the language and common structural/equational basis of Definition~\ref{def:p10-contextual-adjoint-presentations}. All displayed judgments are well formed. \(\Gamma\) is a finite list of distinct sorted variables; \(\Phi\) is a finite hypothesis list. Formula abstractions and capture-avoiding substitution have the conventions of Definition~\ref{def:p10-contextual-first-order}. Down and up in rule names indicate the directions of the source's double rule line; a two-premise bi-rule has two distinct upward directions. These are fifteen directed inferences, derived in the natural-deduction presentation by Theorem~\ref{thm:p10-contextual-adjoint-equivalence}, and usable instead of its logical rules over the same common basis. Source: \cite[Figures~4, 6, and~8, pp.~48, 52, and~55]{Pitts1995Draft}.}
\hbrule{rule:contextual-adjoint-and-down}{\(\&\)-Adj-down}{\Phi\vdash\psi\,[\Gamma]\qquad\Phi\vdash\psi'\,[\Gamma]}{\Phi\vdash\psi\mathbin{\&}\psi'\,[\Gamma]}{The same hypothesis list \(\Phi\) occurs in both premises.}
\hbrule{rule:contextual-adjoint-and-up1}{\(\&\)-Adj-up1}{\Phi\vdash\psi\mathbin{\&}\psi'\,[\Gamma]}{\Phi\vdash\psi\,[\Gamma]}{\none}
\hbrule{rule:contextual-adjoint-and-up2}{\(\&\)-Adj-up2}{\Phi\vdash\psi\mathbin{\&}\psi'\,[\Gamma]}{\Phi\vdash\psi'\,[\Gamma]}{\none}
\hbrule{rule:contextual-adjoint-or-down}{\(\vee\)-Adj-down}{\Phi,\psi\vdash\theta\,[\Gamma]\qquad\Phi,\psi'\vdash\theta\,[\Gamma]}{\Phi,\psi\vee\psi'\vdash\theta\,[\Gamma]}{The same retained hypothesis list \(\Phi\) occurs in both premises and the conclusion.}
\hbrule{rule:contextual-adjoint-or-up1}{\(\vee\)-Adj-up1}{\Phi,\psi\vee\psi'\vdash\theta\,[\Gamma]}{\Phi,\psi\vdash\theta\,[\Gamma]}{\none}
\hbrule{rule:contextual-adjoint-or-up2}{\(\vee\)-Adj-up2}{\Phi,\psi\vee\psi'\vdash\theta\,[\Gamma]}{\Phi,\psi'\vdash\theta\,[\Gamma]}{\none}
\hbrule{rule:contextual-adjoint-imp-down}{\(\Rightarrow\)-Adj-down}{\Phi,\psi\vdash\psi'\,[\Gamma]}{\Phi\vdash\psi\Rightarrow\psi'\,[\Gamma]}{\none}
\hbrule{rule:contextual-adjoint-imp-up}{\(\Rightarrow\)-Adj-up}{\Phi\vdash\psi\Rightarrow\psi'\,[\Gamma]}{\Phi,\psi\vdash\psi'\,[\Gamma]}{\none}
\hbrule{rule:contextual-adjoint-false}{false-Adj}{}{\Phi,\mathsf{false}\vdash\psi\,[\Gamma]}{The source bi-rule has an empty upper side, so there is no substantive upward inference.}
\hbrule{rule:contextual-adjoint-forall-down}{\(\forall\)-Adj-down}{\Phi\vdash\psi(x)\,[\Gamma,x:\sigma]}{\Phi\vdash\forall_\sigma(\psi)\,[\Gamma]}{\(\psi\) is a unary formula abstraction with body in \(\Gamma,x:\sigma\). The variable \(x\) is fresh for \(\Gamma\), and every member of \(\Phi\) is well formed in \(\Gamma\).}
\hbrule{rule:contextual-adjoint-forall-up}{\(\forall\)-Adj-up}{\Phi\vdash\forall_\sigma(\psi)\,[\Gamma]}{\Phi\vdash\psi(x)\,[\Gamma,x:\sigma]}{\(\psi\) is a unary formula abstraction with body in \(\Gamma,x:\sigma\). The variable \(x\) is fresh for \(\Gamma\), and every member of \(\Phi\) is well formed in \(\Gamma\).}
\hbrule{rule:contextual-adjoint-exists-down}{\(\exists\)-Adj-down}{\Phi,\psi(x)\vdash\theta\,[\Gamma,x:\sigma]}{\Phi,\exists_\sigma(\psi)\vdash\theta\,[\Gamma]}{\(\psi\) is a unary formula abstraction with body in \(\Gamma,x:\sigma\). The variable \(x\) is fresh for \(\Gamma\); all members of \(\Phi\) and \(\theta\) are well formed in \(\Gamma\), so none has \(x\) free.}
\hbrule{rule:contextual-adjoint-exists-up}{\(\exists\)-Adj-up}{\Phi,\exists_\sigma(\psi)\vdash\theta\,[\Gamma]}{\Phi,\psi(x)\vdash\theta\,[\Gamma,x:\sigma]}{\(\psi\) is a unary formula abstraction with body in \(\Gamma,x:\sigma\). The variable \(x\) is fresh for \(\Gamma\); all members of \(\Phi\) and \(\theta\) are well formed in \(\Gamma\), so none has \(x\) free.}
\hbrule{rule:contextual-adjoint-equality-down}{\(=\)-Adj-down}{\Phi\vdash\psi(x,x)\,[\Gamma,x:\sigma]}{\Phi,(x=_\sigma x')\vdash\psi(x,x')\,[\Gamma,x:\sigma,x':\sigma]}{The variables \(x,x'\) are distinct and fresh for \(\Gamma\). The binary abstraction \(\psi\) has body in \(\Gamma,x:\sigma,x':\sigma\). Every member of \(\Phi\) is well formed in \(\Gamma,x:\sigma\); it may contain \(x\) free but not \(x'\).}
\hbrule{rule:contextual-adjoint-equality-up}{\(=\)-Adj-up}{\Phi,(x=_\sigma x')\vdash\psi(x,x')\,[\Gamma,x:\sigma,x':\sigma]}{\Phi\vdash\psi(x,x)\,[\Gamma,x:\sigma]}{The variables \(x,x'\) are distinct and fresh for \(\Gamma\). The binary abstraction \(\psi\) has body in \(\Gamma,x:\sigma,x':\sigma\). Every member of \(\Phi\) is well formed in \(\Gamma,x:\sigma\); it may contain \(x\) free but not \(x'\).}
